\begin{table}[tbp]
\centering
\caption{问题三：风格迁移方法比较}
\label{tab:t08_q3_methods}
\footnotesize
\setlength{\tabcolsep}{3.5pt}
\begin{tabular}{llrrrrrrrrr}
\toprule
数据 & 方法 & 约束 & n & 可行率 & 风格位移中位数 & QP 成本中位数 & 新增 DMS & 新增双乙酰 & 新增 4VG & 任一维劣化 \\
\midrule
11 个 Pils/Lager 样本 & 无校正 & — & 11 & 100.00\% & 0.911 & 0.091 & 8 & 6 & 0 & 7/11（63.64\%） \\
11 个 Pils/Lager 样本 & OT 投影 & 松 & 11 & 100.00\% & 0.894 & 0.100 & 0 & 0 & 0 & 7/11（63.64\%） \\
11 个 Pils/Lager 样本 & OT 投影 & 严 & 11 & 100.00\% & 0.889 & 0.110 & 0 & 0 & 0 & 7/11（63.64\%） \\
11 个 Pils/Lager 样本 & OT 条件化 & 松 & 11 & 100.00\% & 0.773 & 0.119 & 0 & 0 & 0 & 7/11（63.64\%） \\
11 个 Pils/Lager 样本 & OT 条件化 & 严 & 11 & 100.00\% & 0.773 & 0.141 & 0 & 0 & 0 & 7/11（63.64\%） \\
11 个 Pils/Lager 样本 & 联合 QP & 松 & 11 & 100.00\% & 0.821 & 0.236 & 0 & 0 & 0 & 0/11（0.00\%） \\
11 个 Pils/Lager 样本 & 联合 QP & 严 & 11 & 100.00\% & 0.800 & 0.272 & 0 & 0 & 0 & 0/11（0.00\%） \\
924 个风格对 & 无校正 & — & 924 & 100.00\% & 0.854 & 0.091 & 166 & 122 & 0 & 868/924（93.94\%） \\
924 个风格对 & OT 投影 & 松 & 924 & 100.00\% & 0.849 & 0.091 & 0 & 0 & 0 & 869/924（94.05\%） \\
924 个风格对 & 联合 QP & 松 & 924 & 100.00\% & 0.775 & 0.165 & 0 & 0 & 0 & 0/924（0.00\%） \\
\bottomrule
\end{tabular}
\end{table}
% 注（写入图注/表注，勿直接编译）：“新增”指从阈值以下升到阈值以上；“任一维劣化”是 stale_hops、DMS、diacetyl、barnyard、metallic 五个感官维度中至少一个变差。风格位移与 QP 成本均为中位数，越小越接近源风格 / 约束越弱。
